\documentclass[10pt]{article}
\usepackage[english]{babel}
\usepackage{../../../../lib/tex/naproche}
\usepackage{amssymb}
\usepackage{mathtools} % for \coloneq

\usepackage{stex-highlighting}
\providebool{emph} % "\newbool{emph}" does not work...
\setbool{emph}{false}
\colorlet{emphcolor}{violet}
\let\oldemph\emph
\renewcommand\emph[1]{\setbool{emph}{true}\ifbool{forthel}{\textcolor{emphcolor}{\itshape#1}}{\oldemph{#1}}\setbool{emph}{false}}
\renewcommand{\varemph}[1]{\ifbool{emph}{\textcolor{emphcolor}{#1}}{\textcolor{black}{#1}}}

\usepackage[right=6cm,left=3cm,bottom=3cm,marginparwidth=5cm]{geometry}

\usepackage{fancyhdr}
\renewcommand{\sectionmark}[1]{\markboth{#1}{}} 
\def\libarchive{}
\pagestyle{fancy}
\fancyhead[L]{\libarchive}
\fancyhead[C]{\nouppercase\leftmark}  % section title
\fancyhead[R]{\thepage}               % page number
\fancyfoot[C]{}                       % No page number in footer

\usepackage[nobottomtitles]{titlesec}
\titlespacing*{\section}{0pt}{30pt}{0pt}
\titlespacing*{\subsection}{0pt}{30pt}{0pt}
\titlespacing*{\subsubsection}{0pt}{30pt}{0pt}

\documentclass[12pt,oneside]{book}

\usepackage[foundations]{../../lib/tex/naproche}
\usepackage{../../lib/tex/libraries}
\usepackage{graphicx}
\usepackage{float}
\usepackage{caption}
\usepackage{footnote}

\makesavenoteenv{tabular} % Make footnotes work in tabular environments


\title{Foundations of Mathematics}
\author{Marcel Schütz}
\date{2022}

\begin{document}
  \maketitle

  \tableofcontents

  \begin{figure}[H]
    \centering
    \fbox{\includegraphics[width=0.9\linewidth]{./dependency-graph/graph.png}}
    \caption*{Interdependencies of the chapters}
  \end{figure}


  \section*{Introduction}

  This is a library providing a foundation of mathematics based on a
  Kelley-Morse like class theory with urelements.
  It introduces common operations on classes like unions or intersections
  (\cref{chapter:classes}) together with detailed proofs of their algebraic
  properties (\cref{chapter:computation-laws-for-classes}), the symmetric
  difference of two classes (\cref{chapter:symmetric-difference}) and the
  notions of ordered pairs and Cartesian products
  (\cref{chapter:pairs-and-products}) as well as proofs of the algebraic
  properties of the latter (\cref{chapter:computation-laws-for-products}).
  Moreover, it provides common operations on maps (\cref{chapter:maps}), various
  properties of images and preimages (\cref{chapter:image-and-preimage}) and the
  notions of injectivity, surjectivity, bijectivity
  (\cref{chapter:injections-surjections-bijections}) and invertibility of maps
  (\cref{chapter:invertible-maps}).
  The library provides an axiom system characterizing sets (\cref{chapter:sets})
  and, furthermore, it covers the notions of binary relations
  (\cref{chapter:binary-relations}), fixed-points of subset preserving maps
  (\cref{chapter:fixed-points}), including and equinumerosity
  (\cref{chapter:equinumerosity}).

  As two famous results it includes the Knaster-Tarski fixed point theorem
  (\cref{FOUNDATIONS_12_8420450166112256}) and the Cantor-Schröder-Bernstein
  theorem (\cref{FOUNDATIONS_13_1913663275401216}).

  \paragraph*{Usage.}
  At the very beginning of each chapter you can find the name of its source
  file, e.g. \path{foundations/sections/01_classes.ftl.tex} for
  \cref{chapter:classes}. This filename can be used to import the chapter via
  \Naproche's \texttt{readtex} instruction to another ForTheL text, e.g.:
  \begin{center}
    \verb`[readtex \path{foundations/sections/01_classes.ftl.tex}]`
  \end{center}

  \paragraph*{Checking times.}
  The checking times for each of the chapters may vary from computer to
  computer, but on mid-range hardware they are likely to be similar to those
  given in table below:

  \begin{center}
    \begin{tabular}{c|c|c}

      & \multicolumn{2}{c}{\textbf{Checking time}}
      \\
      \textbf{Chapter}
      & \textbf{without dependencies}     & \textbf{with dependencies}
      \\ \hline
      \ref{chapter:classes}
      & 00:04 min                         & 00:04 min
      \\
      \ref{chapter:computation-laws-for-classes}
      & 00:12 min                         & 00:16 min
      \\
      \ref{chapter:symmetric-difference}
      & 00:32 min                         & 00:48 min
      \\
      \ref{chapter:pairs-and-products}
      & 00:08 min                         & 00:12 min
      \\
      \ref{chapter:computation-laws-for-products}
      & 01:36 min                         & 01:56 min
      \\
      \ref{chapter:maps}
      & 01:13 min                         & 01:25 min
      \\
      \ref{chapter:image-and-preimage}
      & 01:28 min                         & 02:53 min
      \\
      \ref{chapter:injections-surjections-bijections}
      & 00:38 min                         & 02:03 min
      \\
      \ref{chapter:invertible-maps}
      & 02:20 min                         & 04:23 min
      \\
      \ref{chapter:sets}
      & 02:17 min                         & 06:40 min
      \\
      \ref{chapter:binary-relations}
      & 00:14 min                         & 06:54 min
      \\
      \ref{chapter:fixed-points}
      & 00:33 min                         & 07:13 min
      \\
      \ref{chapter:equinumerosity}
      & 01:48 min                         & 09:01 min
    \end{tabular}
  \end{center}


  \subfile{sections/01_classes.ftl.tex}
  \subfile{sections/02_computation-laws-for-classes.ftl.tex}
  \subfile{sections/03_symmetric-difference.ftl.tex}
  \subfile{sections/04_pairs-and-products.ftl.tex}
  \subfile{sections/05_computation-laws-for-products.ftl.tex}
  \subfile{sections/06_maps.ftl.tex}
  \subfile{sections/07_image-and-preimage.ftl.tex}
  \subfile{sections/08_injections-surjections-bijections.ftl.tex}
  \subfile{sections/09_invertible-maps.ftl.tex}
  \subfile{sections/10_sets.ftl.tex}
  \subfile{sections/11_binary-relations.ftl.tex}
  \subfile{sections/12_fixed-points.ftl.tex}
  \subfile{sections/13_equinumerosity.ftl.tex}
\end{document}

\usepackage{amssymb}

\newcommand{\Nat}{\mathbb{N}}
\newcommand{\Prime}{\mathbb{P}}
\renewcommand{\succ}{\textrm{succ}}
\newcommand{\pred}{\textrm{pred}}
\newcommand{\add}{\textrm{add}}
\newcommand{\mul}{\textrm{mul}}
\renewcommand{\exp}{\textrm{exp}}
\newcommand{\fac}{\textrm{fac}}
\renewcommand{\div}{\mathop{\textrm{div}}}
\renewcommand{\mod}{\mathop{\textrm{mod}}}

\begin{document}
  \begin{imports}
    \begin{forthel}
      %[prove off][check off]
      [readtex \path{libraries/source/arithmetics/05_subtraction.ftl.tex}]
      %[prove on][check on]
    \end{forthel}
  \end{imports}


  \section{Definition of Multiplication}

  \begin{forthel}
    \begin{signature}\printlabel{ARITHMETIC_06_6626346484629504}
      Let $n, m$ be natural numbers.
      $n \cdot m$ is a natural number.
    \end{signature}

    Let the product of $n$ and $m$ stand for $n \cdot m$.
  \end{forthel}

  \begin{forthel}
    \begin{axiom}\printlabel{ARITHMETIC_06_8941041092657152}
      Let $n$ be a natural number.
      Then $n \cdot 0 = 0$.
    \end{axiom}
  \end{forthel}

  \begin{forthel}
    \begin{axiom}\printlabel{ARITHMETIC_06_2211275408932864}
      Let $n, m$ be natural numbers.
      Then $n \cdot (m + 1) = (n \cdot m) + n$.
    \end{axiom}
  \end{forthel}


  \section{Computation Laws}

  \subsection{Distributivity}

  \begin{forthel}
    \begin{proposition}\printlabel{ARITHMETIC_06_9001524774567936}
      Let $n, m, k$ be natural numbers.
      Then \[ n \cdot (m + k) = (n \cdot m) + (n \cdot k). \]
    \end{proposition}
    \begin{proof}
      Define $\Phi = \{ k' \in \Nat \mid n \cdot (m + k') =
      (n \cdot m) + (n \cdot k') \}$.

      (1) $0$ is an element of $\Phi$.
      Indeed $n \cdot (m + 0)
        = n \cdot m
        = (n \cdot m) + 0
        = (n \cdot m) + (n \cdot 0)$.

      (2) For all $k' \in \Phi$ we have $k' + 1 \in \Phi$. \\
      Proof.
        Let $k'\in \Phi$.
        Then
        \[  n \cdot (m + (k' + 1))                  \]
        \[    = n \cdot ((m + k') + 1)              \]
        \[    = (n \cdot (m + k')) + n              \]
        \[    = ((n \cdot m) + (n \cdot k')) + n    \]
        \[    = (n \cdot m) + ((n \cdot k') + n)    \]
        \[    = (n \cdot m) + (n \cdot (k' + 1)).   \]
        Hence $n \cdot (m + (k' + 1)) = (n \cdot m) + (n \cdot (k' + 1))$.
        Thus $k' + 1 \in \Phi$.
      Qed.

      Thus every natural number is contained in $\Phi$.
      Therefore $n \cdot (m + k) = (n \cdot m) + (n \cdot k)$.
    \end{proof}
  \end{forthel}

  \begin{forthel}
    \begin{proposition}\printlabel{ARITHMETIC_06_5742967566368768}
      Let $n, m, k$ be natural numbers.
      Then \[ (n + m) \cdot k = (n \cdot k) + (m \cdot k). \]
    \end{proposition}
    \begin{proof}
      Define $\Phi = \{ k' \in \Nat \mid (n + m) \cdot k' =
      (n \cdot k') + (m \cdot k') \}$.

      (1) $0$ belongs to $\Phi$.
      Indeed $(n + m) \cdot 0
        = 0
        = 0 + 0
        = (n \cdot 0) + (m \cdot 0)$.

      (2) For all $k' \in \Phi$ we have $k' + 1 \in \Phi$. \\
      Proof.
        Let $k' \in \Phi$.
        Then
        \[  (n + m) \cdot (k' + 1)                        \]
        \[    = ((n + m) \cdot k') + (n + m)              \]
        \[    = ((n \cdot k') + (m \cdot k')) + (n + m)   \]
        \[    = (((n \cdot k') + (m \cdot k')) + n) + m   \]
        \[    = ((n \cdot k') + ((m \cdot k') + n)) + m   \]
        \[    = ((n \cdot k') + (n + (m \cdot k'))) + m   \]
        \[    = (((n \cdot k') + n) + (m \cdot k')) + m   \]
        \[    = ((n \cdot k') + n) + ((m \cdot k') + m)   \]
        \[    = (n \cdot (k' + 1)) + (m \cdot (k' + 1)).  \]
        Thus $(n + m) \cdot (k' + 1) = (n \cdot (k' + 1)) + (m \cdot (k' + 1))$.
      Qed.

      Thus every natural number is an element of $\Phi$.
      Therefore $(n + m) \cdot k = (n \cdot k) + (m \cdot k)$.
    \end{proof}
  \end{forthel}


  \subsection{Multiplication with $1$ and $2$}

  \begin{forthel}
    \begin{proposition}\printlabel{ARITHMETIC_06_2910559821365248}
      Let $n$ be a natural number.
      Then \[ n \cdot 1 = n. \]
    \end{proposition}
    \begin{proof}
      $n \cdot 1
        = n \cdot (0 + 1)
        = (n \cdot 0) + n
        = 0 + n
        = n$.
    \end{proof}
  \end{forthel}

  \begin{forthel}
    \begin{corollary}\printlabel{ARITHMETIC_06_5679541582299136}
      Let $n$ be a natural number.
      Then \[ n \cdot 2 = n + n. \]
    \end{corollary}
    \begin{proof}
      $n \cdot 2
        = n \cdot (1 + 1)
        = (n \cdot 1) + n
        = n + n$.
    \end{proof}
  \end{forthel}


  \subsection{Associativity}

  \begin{forthel}
    \begin{proposition}\printlabel{ARITHMETIC_06_347295585402880}
      Let $n, m, k$ be natural numbers.
      Then \[ n \cdot (m \cdot k) = (n \cdot m) \cdot k. \]
    \end{proposition}
    \begin{proof}
      Define $\Phi = \{ k' \in \Nat \mid n \cdot (m \cdot k') =
      (n \cdot m) \cdot k' \}$.

      (1) $0$ is contained in $\Phi$.
      Indeed $n \cdot (m \cdot 0)
        = n \cdot 0
        = 0
        = (n \cdot m) \cdot 0$.

      (2) For all $k' \in \Phi$ we have $k' + 1 \in \Phi$. \\
      Proof.
        Let $k' \in \Phi$.
        Then
        \[  n \cdot (m \cdot (k' + 1))                          \]
        \[    = n \cdot ((m \cdot k') + m)                      \]
        \[    = (n \cdot (m \cdot k')) + (n \cdot m)            \]
        \[    = ((n \cdot m) \cdot k') + (n \cdot m)            \]
        \[    = ((n \cdot m) \cdot k') + ((n \cdot m) \cdot 1)  \]
        \[    = (n \cdot m) \cdot (k' + 1).                     \]
      Qed.

      Hence every natural number is contained in $\Phi$ (by \cref{ARITHMETIC_03_647949900054528}).
      Thus $n \cdot (m \cdot k) = (n \cdot m) \cdot k$.
    \end{proof}
  \end{forthel}


  \subsection{Commutativity}

  \begin{forthel}
    \begin{proposition}\printlabel{ARITHMETIC_06_1764759896588288}
      Let $n, m$ be natural numbers.
      Then \[ n \cdot m = m \cdot n. \]
    \end{proposition}
    \begin{proof}
      Define $\Phi = \{ m' \in \Nat \mid n \cdot m' = m' \cdot n \}$.

      (1) $0$ is contained in $\Phi$. \\
      Proof.
        Define $\Psi = \{ n' \in \Nat \mid n' \cdot 0 = 0 \cdot n' \}$.

        (1a) $0$ is contained in $\Psi$.

        (1b) For all $n' \in \Psi$ we have $n' + 1 \in \Psi$. \\
        Proof.
          Let $n' \in \Psi$.
          Then
          \[ (n' + 1) \cdot 0
            = 0
            = n' \cdot 0
            = 0 \cdot n'
            = (0 \cdot n') + 0
            = 0 \cdot (n' + 1). \]
        Qed.

        Hence every natural number is contained in $\Psi$ (by \cref{ARITHMETIC_03_647949900054528}).
        Thus $n \cdot 0 = 0 \cdot n$.
      Qed.

      (2) $1$ belongs to $\Phi$. \\
      Proof.
        Define $\Theta = \{ n' \in \Nat \mid n' \cdot 1 = 1 \cdot n' \}$.

        (2a) $0$ is contained in $\Theta$.

        (2b) For all $n' \in \Theta$ we have $n' + 1 \in \Theta$. \\
        Proof.
          Let $n' \in \Theta$.
          Then
          \[  (n' + 1) \cdot 1        \]
          \[    = (n' \cdot 1) + 1    \]
          \[    = (1 \cdot n') + 1    \]
          \[    = 1 \cdot (n' + 1).   \]
        Qed.

        Thus every natural number is contained in $\Theta$ (by \cref{ARITHMETIC_03_647949900054528}).
        Therefore $n \cdot 1 = 1 \cdot n$.
      Qed.

      (3) For all $m' \in \Phi$ we have $m' + 1 \in \Phi$. \\
      Proof.
        Let $m' \in \Phi$.
        Then
        \[  n \cdot (m' + 1)                \]
        \[    = (n \cdot m') + (n \cdot 1)  \]
        \[    = (m' \cdot n) + (1 \cdot n)  \]
        \[    = (1 \cdot n) + (m' \cdot n)  \]
        \[    = (1 + m') \cdot n            \]
        \[    = (m' + 1) \cdot n.           \]
        Indeed $((1 \cdot n) + (m' \cdot n)) = (1 + m') \cdot n$. %!
      Qed.

      Hence every natural number is contained in $\Phi$ (by \cref{ARITHMETIC_03_647949900054528}).
      Thus $n \cdot m = m \cdot n$.
    \end{proof}
  \end{forthel}


  \subsection{Non-Existence of Zero-Divisors}

  \begin{forthel}
    \begin{proposition}\printlabel{ARITHMETIC_06_3843962875936768}
      Let $n, m$ be natural numbers such that $n \cdot m = 0$.
      Then $n = 0$ or $m = 0$.
    \end{proposition}
    \begin{proof}
      Suppose $n, m \neq 0$.
      Take natural numbers $n', m'$ such that $n = (n' + 1)$ and $m = (m' + 1)$.
      Then
      \[  0                                     \]
      \[    = n \cdot m                         \]
      \[    = (n' + 1) \cdot (m' + 1)           \]
      \[    = ((n' + 1) \cdot m') + (n' + 1)    \]
      \[    = (((n' + 1) \cdot m') + n') + 1.   \]
      Indeed $(n' + 1) \cdot (m' + 1) = ((n' + 1) \cdot m') + (n' + 1)$.
      Hence $0 = k + 1$ for some natural number $k$.
      Contradiction.
    \end{proof}
  \end{forthel}


  \subsection{Cancellation}

  \begin{forthel}
    \begin{proposition}\printlabel{ARITHMETIC_06_31055184658432}
      Let $n, m, k$ be natural numbers.
      Assume $k \neq 0$.
      Then \[ n \cdot k = m \cdot k \implies n = m. \]
    \end{proposition}
    \begin{proof}
      Define $\Phi = \{ n' \in \Nat \mid$ for all $m' \in \Nat$ if
      $n' \cdot k = m' \cdot k$ and $k \neq 0$ then $n' = m' \}$.

      (1) $0$ is contained in $\Phi$. \\
      Proof.
        Let $m' \in \Nat$.
        Assume $0 \cdot k = m' \cdot k$ and $k \neq 0$.
        Then $m' \cdot k = 0$.
        Hence $m' = 0$ or $k = 0$.
        Thus $m' = 0$.
      Qed.

      (2) For all $n' \in \Phi$ we have $n' + 1 \in \Phi$. \\
      Proof.
        Let $n' \in \Phi$.

        Let us show that for all $m' \in \Nat$ if $(n' + 1) \cdot k =
        m' \cdot k$ and $k \neq 0$ then $n' + 1 = m'$.
          Let $m' \in \Nat$.
          Assume $(n' + 1) \cdot k = m' \cdot k$ and $k \neq 0$.

          Case $m' = 0$.
            Then $(n' + 1) \cdot k = 0$.
            Hence $n' + 1 = 0$.
            Contradiction.
          End.

          Case $m' \neq 0$.
            Take a natural number $l$ such that $m' = l + 1$.
            Then $(n' + 1) \cdot k = (l + 1) \cdot k$.
            Hence $(n' \cdot k) + k
              = (n' \cdot k) + (1 \cdot k)
              = (n' \cdot k) + k
              = (l + 1) \cdot k
              = (l \cdot k) + (1 \cdot k)
              = (l \cdot k) + k$.
            Thus $n' \cdot k = l \cdot k$ (by \cref{ARITHMETIC_03_3137702874578944}).
            Indeed $n' \cdot k$ and $l \cdot k$ are natural numbers.
            Then we have $n' = l$.
            Indeed if $n' \cdot k = l \cdot k$ and $k \neq 0$ then $n' = l$.
            Therefore $n' + 1 = l + 1 = m'$.
          End.
        End.

        Hence $n' + 1 \in \Phi$.
      Qed.

      Thus every natural number is contained in $\Phi$ (by \cref{ARITHMETIC_03_647949900054528}).
      Therefore if $n \cdot k = m \cdot k$ then $n = m$.
    \end{proof}
  \end{forthel}

  \begin{forthel}
    \begin{corollary}\printlabel{ARITHMETIC_06_8575191374364672}
      Let $n, m, k$ be natural numbers.
      Assume $k \neq 0$.
      Then \[ k \cdot n = k \cdot m \implies n = m. \]
    \end{corollary}
    \begin{proof}
      Assume $k \cdot n = k \cdot m$.
      We have $k \cdot n = n \cdot k$ and $k \cdot m = m \cdot k$.
      Hence $n \cdot k = m \cdot k$.
      Thus $n = m$ (by \cref{ARITHMETIC_06_31055184658432}).
    \end{proof}
  \end{forthel}


  \section{Ordering and Multiplication}

  \begin{forthel}
    \begin{proposition}\printlabel{ARITHMETIC_06_8817333933965312}
      Let $n, m, k$ be natural numbers.
      Assume $k \neq 0$.
      Then \[ n < m \iff n \cdot k < m \cdot k. \]
    \end{proposition}
    \begin{proof}
      Case $n \cdot k < m \cdot k$.
        Define $\Phi = \{ n' \in \Nat \mid$ if $n' \cdot k < m \cdot k$ then
        $n' < m \}$.

        (1) $\Phi$ contains $0$.

        (2) For all $n' \in \Phi$ we have $n' + 1 \in \Phi$. \\
        Proof.
          Let $n' \in \Phi$.

          Let us show that if $(n' + 1) \cdot k < m \cdot k$ then $n' + 1 < m$.
            Assume $(n' + 1) \cdot k < m \cdot k$.
            Then $(n' \cdot k) + k < m \cdot k$.
            Hence $n' \cdot k < m \cdot k$.
            Thus $n' < m$.
            Then $n' + 1 \leq m$.
            If $n' + 1 = m$ then $(n' + 1) \cdot k = m \cdot k$.
            Hence $n' + 1 < m$.
          End.
        Qed.

        Therefore every natural number is contained in $\Phi$ (by \cref{ARITHMETIC_03_647949900054528}).
        Consequently $n < m$.
      End.

      Case $n < m$.
        Take a positive natural number $l$ such that $m = n + l$.
        Then $m \cdot k = (n + l) \cdot k = (n \cdot k) + (l \cdot k)$.
        $l \cdot k$ is positive.
        Hence $n \cdot k < m \cdot k$.
      End.
    \end{proof}
  \end{forthel}

  \begin{forthel}
    \begin{corollary}\printlabel{ARITHMETIC_06_5048640368279552}
      Let $n, m, k$ be natural numbers.
      Assume $k \neq 0$.
      Then \[ n < m \iff k \cdot n < k \cdot m \]
      (by \cref{ARITHMETIC_06_8817333933965312}, \cref{ARITHMETIC_06_1764759896588288}).
    \end{corollary}
  \end{forthel}

  \begin{forthel}
    \begin{proposition}\printlabel{ARITHMETIC_06_1826268599287808}
      Let $n, m, k$ be natural numbers.
      Then \[ n, m > k \implies n \cdot m > k. \]
    \end{proposition}
    \begin{proof}
      Define $\Phi = \{ n' \in \Nat \mid$ if $n', m > k$ then $n' \cdot m > k \}$.

      (1) $\Phi$ contains $0$.

      (2) For all $n' \in \Phi$ we have $n' + 1 \in \Phi$. \\
      Proof.
        Let $n' \in \Phi$.

        Let us show that if $n' + 1, m > k$ then $(n' + 1) \cdot m > k$.
          Assume $n' + 1, m > k$.
          Then $(n' + 1) \cdot m = (n' \cdot m) + m$.
          If $n' = 0$ then
          $(n' \cdot m) + m
            = 0 + m
            = m
            > k$.
          If $n' \neq 0$ then
          $(n' \cdot m) + m
            > m
            > k$.
          Indeed if $n' \neq 0$ then $n' \cdot m > 0$.
          Indeed $m > 0$.
          Hence $(n' + 1) \cdot m > k$.
        Qed.
      Qed.

      Thus every natural number is contained in $\Phi$ (by \cref{ARITHMETIC_03_647949900054528}).
      Therefore if $n, m > k$ then $n \cdot m > k$.
    \end{proof}
  \end{forthel}

  \begin{forthel}
    \begin{corollary}\printlabel{ARITHMETIC_06_1751605544222720}
      Let $n, m, k$ be natural numbers.
      Then \[ n \leq m \implies k \cdot n \leq k \cdot m. \]
    \end{corollary}
  \end{forthel}

  \begin{forthel}
    \begin{corollary}\printlabel{ARITHMETIC_06_3965209318260736}
      Let $n, m, k$ be natural numbers.
      Assume $k \neq 0$.
      Then \[ k \cdot n \leq k \cdot m \implies n \leq m. \]
    \end{corollary}
    \begin{proof}
      Assume $k \cdot n \leq k \cdot m$.
      Case $k \cdot n = k \cdot m$. Obvious.
      Case $k \cdot n < k \cdot m$. Obvious.
    \end{proof}
  \end{forthel}

  \begin{forthel}
    \begin{corollary}\printlabel{ARITHMETIC_06_8946886668976128}
      Let $n, m, k$ be natural numbers.
      Then \[ n \leq m \implies n \cdot k \leq m \cdot k. \]
    \end{corollary}
  \end{forthel}

  \begin{forthel}
    \begin{corollary}\printlabel{ARITHMETIC_06_4374428949413888}
      Let $n, m, k$ be natural numbers.
      Assume $k \neq 0$.
      Then \[ n \cdot k \leq m \cdot k \implies n \leq m. \]
    \end{corollary}
  \end{forthel}

  \begin{forthel}
    \begin{proposition}\printlabel{ARITHMETIC_06_8813409145454592}
      Let $n, m, k$ be natural numbers.
      Assume $m > 0$ and $k > 1$.
      Then $k \cdot m > m$.
    \end{proposition}
    \begin{proof}
      Take a natural number $l$ such that $k = l + 2$.
      Then
      \[  k \cdot m                       \]
      \[    = (l + 2) \cdot m             \]
      \[    = (l \cdot m) + (2 \cdot m)   \]
      \[    = (l \cdot m) + (m + m)       \]
      \[    = ((l \cdot m) + m) + m       \]
      \[    = ((l + 1) \cdot m) + m       \]
      \[    \geq 1 + m                    \]
      \[    > m.                          \]
      Indeed $((l + 1) \cdot m) + m \geq 1 + m$.
    \end{proof}
  \end{forthel}


  \section{Multiplication and Subtraction}

  \begin{forthel}
    \begin{proposition}\printlabel{ARITHMETIC_06_5458841930039296}
      Let $n, m, k$ be natural numbers such that $n \geq m$.
      Then \[ (n - m) \cdot k = (n \cdot k) - (m \cdot k). \]
    \end{proposition}
    \begin{proof}
      We have
      \[  ((n - m) \cdot k) + (m \cdot k)                 \]
      \[    = ((n - m) + m) \cdot k                       \]
      \[    = n \cdot k                                   \]
      \[    = ((n \cdot k) - (m \cdot k)) + (m \cdot k).  \]

      Hence $(n - m) \cdot k = (n \cdot k) - (m \cdot k)$.
    \end{proof}
  \end{forthel}

  \begin{forthel}
    \begin{corollary}\printlabel{ARITHMETIC_06_8461123277815808}
      Let $n, m, k$ be natural numbers such that $n \geq m$.
      Then \[ k \cdot (n - m) = (k \cdot n) - (k \cdot m). \]
    \end{corollary}
  \end{forthel}
\end{document}
